\section{Stokes 定理}\label{sec:stokes-theorem}

本节证明一个非常重要的定理.

\begin{RiemannTheorem}[Stokes]\label{thm:stokes}
设 $M$ 是 $n$ 维定向光滑带边界流形, $\omega$ 是 $M$ 上具有紧支集的 $(n-1)$-微分形式, $i:\partial M\to M$ 是边界 $\partial M$ 到 $M$ 的包含映射. 则
\[
  \int_{\partial M}i^*\omega=\int_M\mathrm{d}\omega,
\]
其中 $\partial M$ 取诱导定向 (参见第~\ref{sec:manifolds-with-boundary}~节).
\end{RiemannTheorem}

\begin{RiemannProof}
取从属于 $M$ 的一个开覆盖的单位分解 $\{\rho_i\}_{i\in I}$, 该开覆盖由与定向相容的坐标邻域组成. 则 $\omega=\sum_{i\in I}\rho_i\omega$, 其中可假设 $I$ 是有限的 ($\omega$ 有紧支集), 因而
\[
  \mathrm{d}\omega=\mathrm{d}\sum_{i\in I}\rho_i\omega
  =\sum_{i\in I}\mathrm{d}(\rho_i\omega).
\]
由积分的线性, 有
\[
  \int_M\mathrm{d}\omega=\sum_{i\in I}\int_M\mathrm{d}(\rho_i\omega)
  \quad\text{且}\quad
  \int_{\partial M}i^*\omega=\sum_{i\in I}\int_{\partial M}i^*(\rho_i\omega).
\]
所以, 为证明本定理, 只需考虑 $\operatorname{supp}\omega$ 包含在覆盖中的一个坐标邻域内的情形. 因此, 考虑紧支集包含在一个坐标邻域 $W$ 内的 $(n-1)$-形式 $\omega$. 设 $\varphi:U\to W$ 是相应的参数化, 其中可假设 $U$ 有界 ($\operatorname{supp}(\varphi^*\omega)$ 是紧的). 则 $\omega$ 在 $U$ 上的表示可写成
\[
  \varphi^*\omega=\sum_{j=1}^na_j\,\mathrm{d}x^1\wedge\cdots\wedge
  \mathrm{d}x^{j-1}\wedge\mathrm{d}x^{j+1}\wedge\cdots\wedge\mathrm{d}x^n,
\]
(其中每个 $a_j:U\to\mathbb{R}$ 都是 $C^\infty$ 函数), 且
\[
  \varphi^*\mathrm{d}\omega=\mathrm{d}\varphi^*\omega
  =\sum_{j=1}^n(-1)^{j-1}\frac{\partial a_j}{\partial x^j}
    \,\mathrm{d}x^1\wedge\cdots\wedge\mathrm{d}x^n.
\]
可用下式将函数 $a_j$ 延拓为 $\mathbb{H}^n$ 上的 $C^\infty$ 函数:
\[
  a_j(x^1,\cdots,x^n)=
  \begin{cases}
    a_j(x^1,\cdots,x^n)&\text{若 }(x^1,\ldots,x^n)\in U\\
    0&\text{若 }(x^1,\ldots,x^n)\in\mathbb{H}^n\setminus U.
  \end{cases}
\]
若 $W\cap\partial M=\varnothing$, 则 $i^*\omega=0$. 此外, 在 $\mathbb{H}$ 中取包含 $U$ 的矩形 $I$, 由 $b_j\leq x^j\leq c_j$ ($j=1,\ldots,n$) 定义, 有
\begin{align*}
  \int_M\mathrm{d}\omega
  &=\int_U\left(\sum_{j=1}^n(-1)^{j-1}\frac{\partial a_j}{\partial x^j}\right)
    \mathrm{d}x^1\cdots\mathrm{d}x^n
  =\sum_{j=1}^n(-1)^{j-1}\int_I\frac{\partial a_j}{\partial x^j}\,\mathrm{d}x^1\cdots\mathrm{d}x^n\\
  &=\sum_{j=1}^n(-1)^{j-1}\int_{\mathbb{R}^{n-1}}
    \left(\int_{b_j}^{c_j}\frac{\partial a_j}{\partial x^j}\,\mathrm{d}x^j\right)
    \mathrm{d}x^1\cdots\mathrm{d}x^{j-1}\,\mathrm{d}x^{j+1}\cdots\mathrm{d}x^n\\
  &=\sum_{j=1}^n(-1)^{j-1}\int_{\mathbb{R}^{n-1}}
    \bigl(a_j(x^1,\ldots,x^{j-1},c_j,x^{j+1},\ldots,x^n)\\
  &\qquad{}-a_j(x^1,\ldots,x^{j-1},b_j,x^{j+1},\ldots,x^n)\bigr)
    \,\mathrm{d}x^1\cdots\mathrm{d}x^{j-1}\,\mathrm{d}x^{j+1}\cdots\mathrm{d}x^n=0,
\end{align*}
其中用到了 Fubini 定理 (参见第~\RGref{subsec:notes-stokes-theorem}~节), 微积分基本定理, 以及 $a_j$ 在 $U$ 外为零这一事实. 因此, 在此情形下, $\int_{\partial M}i^*\omega=\int_M\mathrm{d}\omega=0$.

另一方面, 若 $W\cap\partial M\ne\varnothing$, 则取包含 $U$ 的矩形 $I$, 此时由 $b_j\leq x^j\leq c_j$ ($j=1,\ldots,n-1$) 和 $0\leq x^n\leq c_n$ 定义. 与前一种情形一样, 有
\begin{align*}
  \int_M\mathrm{d}\omega
  &=\int_U\left(\sum_{j=1}^n(-1)^{j-1}\frac{\partial a_j}{\partial x^j}\right)
    \mathrm{d}x^1\cdots\mathrm{d}x^n
    =\sum_{j=1}^n(-1)^{j-1}\int_I\frac{\partial a_j}{\partial x^j}\,\mathrm{d}x^1\cdots\mathrm{d}x^n\\
  &=0+(-1)^{n-1}\int_{\mathbb{R}^{n-1}}
    \left(\int_0^{c_n}\frac{\partial a_n}{\partial x^n}\,\mathrm{d}x^n\right)
    \mathrm{d}x^1\cdots\mathrm{d}x^{n-1}\\
  &=(-1)^{n-1}\int_{\mathbb{R}^{n-1}}
    \bigl(a_n(x^1,\ldots,x^{n-1},c_n)-a_n(x^1,\ldots,x^{n-1},0)\bigr)
    \,\mathrm{d}x^1\cdots\mathrm{d}x^{n-1}\\
  &=(-1)^n\int_{\mathbb{R}^{n-1}}a_n(x^1,\ldots,x^{n-1},0)\,\mathrm{d}x^1\ldots\mathrm{d}x^{n-1}.
\end{align*}

为了计算 $\int_{\partial M}i^*\omega$, 需要考虑 $\partial M$ 的一个参数化 $\widetilde\varphi$, 它定义在 $\mathbb{R}^{n-1}$ 的一个开子集上, 并在 $\partial M$ 取诱导定向时保持 $\mathbb{R}^{n-1}$ 的标准定向. 为此, 例如可取集合
\[
  \widetilde U=\{(x^1,\ldots,x^{n-1})\in\mathbb{R}^{n-1}\mid
  ((-1)^nx^1,x^2,\ldots,x^{n-1},0)\in U\}
\]
和由下式给出的参数化 $\widetilde\varphi:\widetilde U\to\partial M$:
\[
  \widetilde\varphi(x^1,\ldots,x^{n-1})
    :=\varphi\bigl((-1)^nx^1,x^2,\ldots,x^{n-1},0\bigr).
\]
回顾仅去掉 $\varphi$ 的最后一个坐标所得的 $\partial M$ 上的定向, 是 $\partial M$ 的诱导定向的 $(-1)^n$ 倍 (参见第~\ref{sec:manifolds-with-boundary}~节). 因此, $\widetilde\varphi$ 给出正确的定向. $i:\partial M\to M$ 在这些坐标下的局部表达式 ($\widehat i:\widetilde U\to U$, 且 $\widehat i=\varphi^{-1}\circ i\circ\widetilde\varphi$) 为
\[
  \widehat i(x^1,\ldots,x^{n-1})
    =\bigl((-1)^nx^1,x^2,\ldots,x^{n-1},0\bigr).
\]
因此,
\[
  \int_{\partial M}i^*\omega
  =\int_{\widetilde U}\widetilde\varphi^*i^*\omega
  =\int_{\widetilde U}(i\circ\widetilde\varphi)^*\omega
  =\int_{\widetilde U}(\varphi\circ\widehat i)^*\omega
  =\int_{\widetilde U}\widehat i^*\varphi^*\omega.
\]
此外,
\begin{align*}
  \widehat i^*\varphi^*\omega
  &=\widehat i^*\sum_{j=1}^na_j\,\mathrm{d}x^1\wedge\cdots\wedge
    \mathrm{d}x^{j-1}\wedge\mathrm{d}x^{j+1}\wedge\cdots\wedge\mathrm{d}x^n\\
  &=\sum_{j=1}^n(a_j\circ\widehat i)\,\mathrm{d}\widehat i^1\wedge\cdots\wedge
    \mathrm{d}\widehat i^{j-1}\wedge\mathrm{d}\widehat i^{j+1}\wedge\cdots\wedge\mathrm{d}\widehat i^n\\
  &=(-1)^n(a_n\circ\widehat i)\,\mathrm{d}x^1\wedge\cdots\wedge\mathrm{d}x^{n-1},
\end{align*}
因为 $\mathrm{d}\widehat i^1=(-1)^n\,\mathrm{d}x^1$, $\mathrm{d}\widehat i^n=0$, 且对 $j\ne1$ 和 $j\ne n$, $\mathrm{d}\widehat i^j=\mathrm{d}x^j$. 因而
\begin{align*}
  \int_{\partial M}i^*\omega
  &=(-1)^n\int_{\widetilde U}(a_n\circ\widehat i)\,\mathrm{d}x^1\cdots\mathrm{d}x^{n-1}\\
  &=(-1)^n\int_{\widetilde U}a_n\bigl((-1)^nx^1,x^2,\ldots,x^{n-1},0\bigr)
    \,\mathrm{d}x^1\cdots\mathrm{d}x^{n-1}\\
  &=(-1)^n\int_{\mathbb{R}^{n-1}}a_n(x^1,x^2,\ldots,x^{n-1},0)
    \,\mathrm{d}x^1\cdots\mathrm{d}x^{n-1}=\int_M\mathrm{d}\omega
\end{align*}
(这里用到了换元定理).
\end{RiemannProof}

\begin{RiemannRemark}\label{rem:stokes-on-manifold-without-boundary}
若 $M$ 是 $n$ 维定向可微流形 (即边界 $\partial M=\varnothing$ 的带边界流形), 则由 Stokes 定理的证明可知, 对 $M$ 上任意具有紧支集的 $(n-1)$-微分形式 $\omega$, 有
\[
  \int_M\mathrm{d}\omega=0.
\]
若将空集上的积分定义为零, 这可视为 Stokes 定理的一个特例.
\end{RiemannRemark}

\begin{RiemannExercise}\label{ex:stokes-theorem}
\begin{enumerate}[label=(\arabic*),ref=\arabic*,leftmargin=*]
\item\label{ex:torus-complement-integral-by-stokes} 利用 Stokes 定理验证 Exercise~\ref{ex:integration-on-manifolds}(\ref{ex:integral-on-complement-of-torus-in-sphere}) 的结果.
\item\label{ex:homotopy-invariance-of-integral} (\emph{积分的同伦不变性}) 回顾: 若存在可微映射 $H:\mathbb{R}\times M\to N$, 使得 $H(0,p)=f_0(p)$ 和 $H(1,p)=f_1(p)$, 则称两个映射 $f_0,f_1:M\to N$ 光滑同伦 [参见 Exercise~\ref{ex:differential-forms}(\ref{ex:homotopy-operator-and-poincare-lemma})]. 若 $M$ 是 $n$ 维紧定向流形, $\omega$ 是 $N$ 上的闭 $n$-形式, 证明
\[
  \int_Mf_0^*\omega=\int_Mf_1^*\omega.
\]
\item\label{ex:sphere-vector-field-and-antipodal-homotopy}
\begin{enumerate}[label=(\alph*),leftmargin=*]
\item 设 $X\in\mathfrak{X}(S^n)$ 是处处不为零的向量场. 证明
\[
  H(t,p)=\cos(\pi t)p+\sin(\pi t)\frac{X_p}{\lVert X_p\rVert}
\]
是恒等映射与对径映射之间的光滑同伦, 其中用到了等同关系
\[
  X_p\in T_pS^n\subset T_p\mathbb{R}^{n+1}\cong\mathbb{R}^{n+1}.
\]
\item 利用 Stokes 定理证明
\[
  \int_{S^n}\omega>0,
\]
其中
\[
  \omega=\sum_{i=1}^{n+1}(-1)^{i+1}x^i\,\mathrm{d}x^1\wedge\cdots\wedge
  \mathrm{d}x^{i-1}\wedge\mathrm{d}x^{i+1}\wedge\cdots\wedge\mathrm{d}x^{n+1},
\]
且 $S^n=\partial\{x\in\mathbb{R}^{n+1}\mid\lVert x\rVert\leq1\}$ 取由 $\mathbb{R}^{n+1}$ 的标准定向诱导的定向.
\item 证明若 $n$ 为偶数, 则 $X$ 不存在. 若 $n$ 为奇数, 情况如何?
\end{enumerate}
\item\label{ex:degree-of-map} (\emph{映射的度}) 设 $M,N$ 是 $n$ 维紧连通定向流形, $f:M\to N$ 是光滑映射. 可以证明, 存在实数 $\deg(f)$ (称为 $f$ 的\textbf{度}), 使得对任何 $n$-形式 $\omega\in\Omega^n(N)$,
\[
  \int_Mf^*\omega=\deg(f)\int_N\omega.
\]
\begin{enumerate}[label=(\alph*),leftmargin=*]
\item 证明若 $f$ 不是满射, 则 $\deg(f)=0$.
\item 证明若 $f$ 是保持定向的微分同胚, 则 $\deg(f)=1$; 若 $f$ 是反转定向的微分同胚, 则 $\deg(f)=-1$.
\item 设 $f:M\to N$ 是满射, $q\in N$ 是 $f$ 的一个正则值. 证明 $f^{-1}(q)$ 是有限集, 且 $q$ 在 $N$ 中有一个邻域 $W$, 使得 $f^{-1}(W)$ 是 $M$ 中开集 $V_i$ 的不交并, 而 $f|_{V_i}:V_i\to W$ 是微分同胚.
\item 承认 $f$ 存在正则值, 证明 $\deg(f)$ 是整数.
{\small (\textbf{Remark:} \textbf{Sard 定理}保证同维流形之间可微映射 $f$ 的临界值集合测度为零, 由此保证 $f$ 存在正则值).}
\item 给定 $n\in\mathbb{N}$, 给出一个度为 $n$ 的光滑映射 $f:S^1\to S^1$.
\item 证明同伦的映射有相同的度.
\item 若 $n$ 为偶数, 设 $f:S^n\to S^n$ 是保持定向的微分同胚; 若 $n$ 为奇数, 则设它是反转定向的微分同胚. 证明 $f$ 有不动点, 即有一点 $p\in S^n$ 使得 $f(p)=p$.
{\small (\textbf{提示:} 证明若 $f$ 没有不动点, 则可以构造 $f$ 与对径映射之间的同伦).}
\end{enumerate}
\end{enumerate}
\end{RiemannExercise}
